\chapter{Primes}\label{ch:primes}

\begin{center}\itshape
There is always a prime between $n$ and $2n$. Chebyshev proved it in 1852;\\
Erd\H{o}s, at nineteen, found the proof given here, which uses nothing but the binomial coefficient $\binom{2n}n$.
\end{center}


Suppose someone hands you a complete list of primes. Multiply its entries and add $1$. Every prime on the list divides the product, so none divides the new number. Yet every integer greater than $1$ has a prime factor. The list cannot have been complete.

This argument proves a surprising statement without predicting a single next prime. If the list begins $2,3,5$, the new number is $31$, which happens to be prime. With a longer list, the number formed this way can be composite. Euclid's proof needs a \emph{new prime factor}, not a new prime value of a particular formula. That is the distinction between a theorem and a tempting but stronger guess.

The second question is quantitative. Infinitely many primes could still be so sparse that large stretches of the number line contain none. Bertrand's theorem rules out one kind of emptiness: for every $n>1$, some prime lies strictly between $n$ and $2n$. The proof in this chapter looks at the binomial coefficient $\binom{2n}{n}$. Its size, and the prime factors that must contribute to it, pull in opposite directions until a prime in the desired interval is forced.

The pictures of primes will help us notice gaps and clusters. They cannot settle a claim about every $n$. We will use exact divisibility for infinitude and inequalities for gaps. The same primes answer both questions, but the two proofs have different jobs.


\portrait{chebyshev}{Pafnuty Chebyshev (1821--1894). Photograph, c.\,1890. Public domain, via Wikimedia Commons.}

\section{Infinitely many primes}

% t1-pass2 motive 6cef39db
Here is the whole idea of Euclid's proof: any proposed finite list of primes misses at least one. Multiply the listed primes and add $1$. A prime factor of the result cannot be on the list. Notice the wording: the new \emph{factor} is prime; the whole product plus $1$ might not be. Once that distinction is secure, we can ask a harder question: not just whether there are infinitely many primes, but how quickly they appear.

\begin{theorem}[Euclid]
There are infinitely many primes.
\end{theorem}
\begin{proof}
Suppose we have a finite list of primes $p_1,\ldots,p_k$.
Form the positive integer
\[
N=p_1p_2\cdots p_k+1.
\]
Since $N>1$, the fundamental theorem of arithmetic
gives at least one prime divisor $q$ of $N$.
If $q$ were one of the listed primes, it would
divide the product $p_1\cdots p_k=N-1$ as well.
A common divisor of $N$ and $N-1$ divides their
difference $1$, but a prime cannot divide $1$.
Therefore $q$ was not on the list. Every finite
list misses a prime, so there are infinitely many.
We did not need $N$ itself to be prime.
\end{proof}

% t1-pass2 why
\paragraph{Why this works.}
What does the product-plus-one construction really buy us? The number $N=p_1\cdots p_k+1$ has at least one prime factor, because it is greater than $1$. None of the listed primes divides $N$: each divides the product $N-1$, so dividing both would force it to divide $1$. The new factor proves the list incomplete. We do not need the stronger, usually false, claim that the whole number $N$ is prime.

\subsection*{What Euclid's argument does, and does not, say}
Read the conclusion of Euclid's proof carefully. Given a finite list $p_1,\ldots,p_k$, put $M=p_1\cdots p_k+1$. The fundamental theorem gives a prime factor $q$ of $M$. If $q$ were on the list, it would divide both $M$ and $M-1$, hence $1$, which is impossible. We have found a new prime \emph{factor}; we have not proved that $M$ itself is prime. In fact $30031=59\cdot509$ is the product of two primes even though it is the first six primes multiplied together, plus $1$.

The number $M$ need not itself be prime. For the first six primes,
\[
2\cdot3\cdot5\cdot7\cdot11\cdot13+1
 =30\,031=59\cdot509.
\]
Both factors are new, which is all the proof needs. Euclid promises an escape from every finite list, not a machine that prints only primes. This distinction is a good habit: prove precisely the claim you need, and do not accidentally order a larger theorem.

\subsection*{A sieve one can hold in the mind}
Write the integers from $2$ to $50$ in a grid. Keep $2$ and cross out all its larger multiples. Keep the next uncrossed number, $3$, and repeat. Continue with $5$ and $7$. At this point you may stop: any composite integer at most $50$ has a prime factor at most $\sqrt{50}$, and we have checked every possible prime up to that bound. This is the Sieve of Eratosthenes. The stopping rule is a theorem about factors, not just a convenience for drawing the picture.

Why stop at $\sqrt N$? If a composite $m\le N$ is $uv$ with $1<u\le v$, then $u\le\sqrt m\le\sqrt N$; a prime factor of $u$ would already have removed $m$. This is a proof of the stopping rule, not just an observation about a few examples.

\begin{figure}[ht]
\centering\fig{ch02-sieve}
\caption{The sieve up to $50$. A cross marks a multiple of $2$ greater than $2$; a slash marks a still unremoved multiple of $3$; a small square marks $25$ and $35$, removed at the $5$-stage; a diamond marks $49$, removed at the $7$-stage. Unmarked entries are primes. Shapes distinguish the stages even in a black-and-white paperback.}
\label{fig:sieve}
\end{figure}

\subsection*{Testing a number and factoring a number}
Testing whether $N$ is prime is different from finding its full prime factorisation. For the first task, divide by primes no larger than $\sqrt N$. If none divides $N$, then $N$ is prime: otherwise one of the two factors of a composite number would be small enough to detect. If a divisor appears, you know only that $N$ is composite. To factor it completely, continue the search on the remaining quotient. Keep this distinction visible when interpreting an algorithm's running time.

For example, $221$ survives division by $2,3,5,7,11$, but $13$ divides it:
$221=13\cdot17$. For $223$, the same tests suffice since $\sqrt{223}<15$;
none succeeds, so $223$ is prime. The sieve instead shares work among all
integers up to a bound. Choosing a method starts with choosing the task.

Counting divisions is not the same as counting input digits. An $s$-digit
integer can be close to $10^s$, so trial division up to its square root
can require testing a range of size comparable to $10^{s/2}$. This explains
why the elementary certificate is useful for small examples but inefficient
for very large inputs. Faster primality tests belong to later study; no
claim about their cost is needed here.

Euler went much further: the primes are numerous enough for their reciprocals to diverge, whereas $\sum\frac1{n^2}$ converges. So primes are, in a precise sense, denser than squares.

\begin{think}
Before reading on: do you expect $\sum_p\frac1p$ to converge or diverge? The squares are sparse enough for $\sum\frac1{n^2}<\infty$, and among $1,\dots,100$ there are $10$ squares but $25$ primes. Euler's idea is to compare $\prod_{p\le x}(1-\frac1p)^{-1}$ with the harmonic sum $\sum_{n\le x}\frac1n$. Can you see why these two are related before looking at the proof?
\end{think}

\begin{theorem}[Euler]\label{thm:euler}
For $x\ge2$, $\displaystyle\sum_{p\le x}\frac1p\ge\ln\ln x-1$. In particular $\sum_p\frac1p=+\infty$.
\end{theorem}
\begin{proof}
Every $n\le x$ factors into primes $\le x$, so expanding the finite product of geometric series,
\[
\begin{aligned}\prod_{p\le x}\left(1-\frac1p\right)^{-1}&=\prod_{p\le x}\sum_{k\ge0}\frac1{p^k}\ge\sum_{n\le x}\frac1n\eqbr\ge\int_1^{\lfloor x\rfloor+1}\frac{\mathrm dt}t\ge\ln x.\end{aligned}
\]
For $0\le t<1$, $-\ln(1-t)-t=\int_0^t\frac u{1-u}\,\mathrm du\le\frac1{1-t}\int_0^t u\,\mathrm du\le\frac{t^2}{1-t}$. With $t=\frac1p$ the last term is $\frac1{p(p-1)}$, and $\sum_p\frac1{p(p-1)}\le\sum_{n\ge2}\frac1{n(n-1)}=1$. Taking logarithms,
$\ln\ln x\le\sum_{p\le x}-\ln\left(1-\frac1p\right)\le\sum_{p\le x}\frac1p+1$.
\end{proof}

% t1-pass2 why
\paragraph{Why this works.}
The Euler product is a bridge between integers and primes. Up to a fixed bound, every integer appears in the expanded product because every integer has a unique prime factorisation. If the sum of prime reciprocals converged, the product over primes would stay bounded. But the harmonic sums on the integer side grow at least like $\ln x$. The comparison forces divergence. Notice which steps involve finite products and where a limit is taken; mixing those two is the tempting shortcut to avoid.

% t1-pass2 examples 6cef39db
\begin{examplex}[Euclid's number need not be prime]
The product of the first six primes plus one is $2\cdot 3\cdot 5\cdot 7\cdot 11\cdot 13+1=30031=59\cdot 509$. Both factors are primes absent from the original list. The proof asked only for one new prime factor.
\end{examplex}

\begin{examplex}[a partial sum of reciprocals]
The primes up to $10$ contribute $1/2+1/3+1/5+1/7=1.176\ldots$, while $\ln\ln 10-1\approx -0.17$ is a weak lower bound and still negative. The inequality of Theorem~\ref{thm:euler} is true and not sharp. By $x=100$ the prime reciprocal sum is already about $1.80$, and $\ln\ln 100-1\approx 0.53$. The divergence is slow, which is why Volume~6's asymptotic is a finer instrument.
\end{examplex}

\section{Chebyshev's bounds}

% t1-pass2 motive 29674c3a
If you count the primes at most $x$, call the answer $\pi(x)$. Euclid says this count grows without bound, but supplies no scale. Compare $\pi(100)$ and $\pi(1000)$ and you will want something more informative than a guarantee that it grows. Chebyshev proves that the count has the rough size $x/\ln x$, up to positive constant factors. The tool is unexpectedly familiar: compare the size of a binomial coefficient with the prime powers in its factorisation. The prime number theorem will sharpen this estimate in Volume~6.

The quantitative sections use the logarithm, integral comparison and the
mean-value tools of Volume~1. The elementary infinitude argument above does
not depend on those analytic tools. Here they measure how many primes there
are, rather than merely showing that the list cannot end.


Let $\pi(x)$ be the number of primes $\le x$. The prime number theorem, proved in Volume~6, says $\pi(x)\sim\frac x{\ln x}$. Chebyshev's elementary bounds already give the right order of magnitude. Everything rests on the central binomial coefficient. Figure~\ref{fig:primecount} shows how close the two are already for small $x$, and how slowly the ratio creeps towards $1$.

\begin{figure}[ht]
\centering
\fig{ch02-primecount}
\caption{Left: the staircase $\pi(x)$ for $x\le100$ against $\frac x{\ln x}$ and the logarithmic integral $\mathrm{Li}(x)=\int_2^x\frac{\mathrm dt}{\ln t}$. Right: the ratios $\pi(x)\big/\frac x{\ln x}$ and $\pi(x)/\mathrm{Li}(x)$ up to $x=10^5$. At $10^5$ the first is still $1.104$, while the second is already $0.996$: $\mathrm{Li}(x)$ is by far the better approximation.}
\label{fig:primecount}
\end{figure}

\begin{think}
Which primes divide $\binom{20}{10}=184\,756$? Factor it by hand, and compare with the list of primes up to $20$. Which primes are missing, and can you explain why $7$ in particular is absent? Your answer is the heart of Erd\H{o}s' proof below.
\end{think}

\begin{lemma}\label{lem:binom}
For $n\ge1$:
\begin{enumerate}[label=(\roman*)]
\item $\dfrac{4^n}{2n}\le\dbinom{2n}n\le4^n$, and $\dbinom{2n+1}n\le4^n$;
\item every prime power $p^a$ dividing $\binom{2n}n$ exactly satisfies $p^a\le2n$;
\item primes $p$ with $\sqrt{2n}<p$ divide $\binom{2n}n$ at most once, and primes with $\frac{2n}3<p\le n$ do not divide it at all, if $n\ge5$.
\end{enumerate}
\end{lemma}
\begin{proof}
(i) $\binom{2n}n$ is the largest of the $2n+1$ terms of $\sum_k\binom{2n}k=4^n$, and is at least the sum $2$ of the two end terms; so $4^n\le(2n-1)\binom{2n}n+2\le2n\binom{2n}n$. Also $\binom{2n+1}n=\binom{2n+1}{n+1}$ and both appear in $\sum_k\binom{2n+1}k=2\cdot4^n$.

(ii) By Legendre (Chapter~\ref{ch:divisibility}),
\[v_p\binom{2n}n=\sum_{k\ge1}\left(\left\lfloor\frac{2n}{p^k}\right\rfloor-2\left\lfloor\frac n{p^k}\right\rfloor\right).\]
Each term is $0$ or $1$ (as $\lfloor2y\rfloor-2\lfloor y\rfloor\in\{0,1\}$), and is $0$ when $p^k>2n$. So $v_p\le\max\{k:p^k\le2n\}$.

(iii) If $p>\sqrt{2n}$ only the term $k=1$ can be non-zero. If $\frac{2n}3<p\le n$ then $2\le\frac{2n}p<3$ and $1\le\frac np<\frac32$, so the $k=1$ term is $2-2=0$, and $p^2>2n$ when $n\ge5$ (as then $\frac{2n}3>\sqrt{2n}$).
\end{proof}

\begin{theorem}[Chebyshev-type bounds]\label{thm:cheb}
\begin{enumerate}[label=(\roman*)]
\item For every real $x\ge1$, $\displaystyle\prod_{p\le x}p<4^x$.
\item For every $n\ge1$, $\pi(2n)\ge\dfrac{n\ln4-\ln(2n)}{\ln(2n)}$. In particular $\pi(x)\ge\frac{\ln2}2\cdot\frac x{\ln x}$ for large $x$.
\end{enumerate}
\end{theorem}
\begin{proof}
(i) It suffices to treat integers $x=N$, by induction. $N=1,2$ are clear, and for even $N>2$ the product equals the one for $N-1$. For $N=2m+1$, every prime in $]m+1,2m+1]$ divides $\binom{2m+1}m=\frac{(2m+1)!}{m!\,(m+1)!}$ (it divides the numerator, not the denominator), so their product is at most $\binom{2m+1}m\le4^m$, and
$\prod_{p\le2m+1}p\le4^m\prod_{p\le m+1}p<4^m\cdot4^{m+1}=4^{2m+1}$.

(ii) Only primes $p\le2n$ divide $\binom{2n}n$, each to a power at most $2n$ by the lemma, so $\frac{4^n}{2n}\le\binom{2n}n\le(2n)^{\pi(2n)}$. Take logarithms.
\end{proof}

% t1-pass2 why
\paragraph{Why this works.}
One bound concerns a product of primes, another the central binomial coefficient. Taking logarithms of a product turns it into a sum, so prime factors become easier to count. The number $\binom{2n}{n}$ is large simply because the $2n$th row of Pascal's triangle has sum $4^n$ and its central entry is largest. Its prime factorisation cannot be too small to match that size. Comparing those descriptions is the heart of the estimates; the constants come from making both bounds explicit.

\subsection*{Where the upper bound for the prime count comes from}
Part~(i) is a bound on a product, rather than directly on a count. Taking logarithms turns it into one:
\[
\vartheta(x):=\sum_{p\le x}\ln p<x\ln4.
\]
For $x>1$, each prime $\sqrt x<p\le x$ contributes more than $\tfrac12\ln x$. There are $\pi(x)-\pi(\sqrt x)$ such primes, so
\[
\begin{aligned}
\pi(x)-\pi(\sqrt x)&\le\frac{2\vartheta(x)}{\ln x},\\
\pi(x)&<\sqrt x+2\ln4\,\frac x{\ln x}.
\end{aligned}
\]
Here $\pi(\sqrt x)\le\sqrt x$ simply counts the integers available. In particular $\pi(x)\le(1+2\ln4)x/\ln x$ for all sufficiently large $x$, because $\ln x\le\sqrt x$ eventually.

For completeness, the lower bound in part~(ii) also extends from even integers to real $x$. Put $n=\lfloor x/2\rfloor$, so $2n\le x$ and $n\ge x/2-1$. Its numerator is positive for large $x$, and then
\[
\pi(x)\ge\pi(2n)\ge
\frac{(x/2-1)\ln4-\ln x}{\ln x}.
\]
For large $x$ the numerator is at least $(\ln2)x/2$. Thus positive constant multiples of $x/\ln x$ bound $\pi(x)$ from both sides. These constants are deliberately modest: this argument discovers the scale of the landscape before measuring every hill.

% t1-pass2 examples 29674c3a
\begin{examplex}[the product of the small primes]
The primes up to $11$ multiply to $2310$, and $4^{11}=4194304$. Theorem~\ref{thm:cheb}(i) says the product up to $x$ is less than $4^x$, and $2310<4^{11}$ is the inequality at a single integer, not the induction. Taking logarithms gives $\theta(11)=\sum_{p\le 11}\ln p=\ln 2310\approx 7.745$, while $11\ln 4\approx 15.25$.
\end{examplex}

\begin{examplex}[a lower bound one can see]
For $n=6$, $\binom{12}{6}=924$ and $4^6/12=341.3$, so $924\ge 4^6/(2\cdot 6)$. The primes up to $12$ are the only ones that can divide this binomial coefficient, and taking logarithms forces several of them to exist. This is the mechanism of the lower bound, before the general estimate.
\end{examplex}

\section{Bertrand's postulate}

% t1-pass2 motive 07e471c9
Try the intervals $(10,20)$ and $(100,200)$. Each contains primes, but a few examples cannot promise this for every starting point. Bertrand's theorem says that for every integer $n>1$ there is a prime \emph{strictly} between $n$ and $2n$. To prove it, suppose the interval were empty. Then account for the prime factors of $\binom{2n}{n}$ and compare that upper bound with a lower bound for the same integer. The contradiction is a contest between two exact estimates.

\begin{think}
Check by hand that there is a prime in $]n,2n]$ for $n=1,\dots,12$. Then imagine there were none for some large $n$: by the last thinking prompt, $\binom{2n}n$ would be made only of primes $\le\frac{2n}3$. Why should that make $\binom{2n}n$ too small?
\end{think}

\begin{theorem}[Bertrand's postulate; Chebyshev 1852, proof by Erd\H{o}s 1932]\label{thm:bertrand}
For every $n\ge1$ there is a prime $p$ with $n<p\le2n$.
\end{theorem}
\begin{proof}
\emph{Small $n$.} The primes $2$, $3$, $5$, $7$, $13$, $23$, $43$, $83$, $163$, $317$, $631$ each are less than twice the previous one. For $1\le n<631$, let $q$ be the largest of them with $q\le n$ (or $q=1$ if $n=1$) and $q'$ the next one: then $n<q'\le2q\le2n$ (equality $q'=2q$ only for $n=1$, where $q'=2$). So the theorem holds for $n\le630$.

\emph{Large $n$.} Let $n\ge631$ and suppose there is no prime in $]n,2n]$. By Lemma~\ref{lem:binom}, the primes dividing $\binom{2n}n$ are all $\le\frac{2n}3$. Those $\le\sqrt{2n}$ (fewer than $\sqrt{2n}$ of them) contribute at most $2n$ each; the others contribute at most $p$, and their product is at most $\prod_{p\le2n/3}p<4^{2n/3}$ by Theorem~\ref{thm:cheb}. Hence
\[
\begin{gathered}\frac{4^n}{2n}\le\binom{2n}n<(2n)^{\sqrt{2n}}\,4^{2n/3},\qq\text{so}\qq h(n)=\frac n3\ln4-\left(1+\sqrt{2n}\right)\ln(2n)<0.\end{gathered}
\]
But $h(n)>0$ for all real $n\ge512$. At $n=512$, $2n=2^{10}$ and $\sqrt{2n}=32$, so $h(512)=\frac{1024}3\ln2-33\cdot10\ln2=\frac{34}3\ln2>0$. And $h$ increases on $[512,\infty[$:
\[
h'(n)=\frac23\ln2-\frac{\ln(2n)}{\sqrt{2n}}-\frac{1+\sqrt{2n}}n,
\]
where $\frac{\ln u}{\sqrt u}$ decreases for $u\ge e^2$, so for $2n\ge1024$ it is at most $\frac{10\ln2}{32}<0.22$, and $\frac{1+\sqrt{2n}}n\le\frac1{512}+\frac{\sqrt2}{\sqrt{512}}<0.07$; while $\frac23\ln2>0.46$. This contradiction proves the theorem.
\end{proof}

% t1-pass2 why
\paragraph{Why this works.}
Assume there is no prime in $(n,2n)$ and examine $\binom{2n}{n}$. Then its prime factors all lie in ranges below $n$. Very small primes may occur repeatedly, and medium primes are controlled by the product estimate. When you multiply those bounds, the result is eventually smaller than the lower bound for the same binomial coefficient. Contradiction. The final finite range still needs checking: an asymptotic inequality cannot quietly cover the small cases.

\subsection*{Reading the proof as a contest of sizes}
The binomial coefficient $\binom{2n}n$ counts the ways to choose a team of $n$ people from $2n$ candidates. It is enormous: at least $4^n/(2n)$. If there were no prime between $n$ and $2n$, all its prime factors would have to come from smaller primes. Primes in $(2n/3,n]$ cancel between numerator and denominator; primes above $\sqrt{2n}$ appear at most once; each remaining prime power is at most $2n$. The resulting upper bound is too small. A counting problem about teams has forced a new prime to exist.

The calculus at the end is only a way to finish this comparison uniformly. Its two jobs are separate: check one starting value of $h$, then prove $h$ increases forever afterwards. Checking many numerical values of $h$ would not establish the second job.

\begin{reminder}[Three different meanings of ``many primes'']
Infinitude means no finite list contains all primes. Divergence of $\sum_p1/p$ says more: their reciprocals cannot have a finite total. Bounds of order $x/\ln x$ say more again, locating the scale of the count up to $x$. None of these statements says that a particular number of the form $M+1$ is prime, or that a short fixed interval must always contain a prime.
\end{reminder}

\subsection*{Two lanes modulo $4$}
Except for $2$, every prime is congruent to $1$ or $3$ modulo $4$. Put the primes on two tracks according to those remainders. The two tracks look irregular; neither one is a simple repeating sequence. Yet there are infinitely many primes in each class. The familiar proof for primes $3$ modulo $4$ modifies Euclid's product, while the class $1$ modulo $4$ will become especially important when we study square roots of $-1$.

\begin{figure}[ht]
\centering
\fig{ch02-mod4-line}
\caption{Primes up to $60$. Circles mark primes congruent to $1$ modulo $4$, squares mark primes congruent to $3$ modulo $4$, and the dot at $2$ is the even prime, which lies in neither odd class.}
\label{fig:mod4-line}
\end{figure}

% t1-pass2 examples 07e471c9
\begin{examplex}[the postulate up to $20$]
For each $n\le 20$ the interval $(n,2n]$ contains a prime: $2\in(1,2]$, $3\in(2,4]$, $5\in(3,6]$, $7\in(4,8]$, $11\in(6,12]$, $13\in(7,14]$, $19\in(10,20]$, $23\in(12,24]$. The general proof reduces the large case to the size of $\binom{2n}{n}$ and checks a finite range by a chain of primes, each less than twice the previous.
\end{examplex}

\begin{examplex}[where the primes in the middle cancel]
In $\binom{20}{10}=184756$, the prime $13$ satisfies $20/3<13\le 10$? No: $13>10$, so $13$ lies in $(10,20]$ and divides the binomial coefficient. The prime $7$ satisfies $20/3<7\le 10$, so it cancels between numerator and denominator. That cancellation is why a gap $(n,2n]$ would leave the binomial coefficient with only small prime factors, and too small.
\end{examplex}

\begin{beauty}
Where this leads. Theorem~\ref{thm:euler} is the input to Mertens's theorems in Volume~6, and Theorem~\ref{thm:cheb} is the Chebyshev bound quoted there before the prime number theorem. Bertrand's postulate is the elementary shadow of the short-interval results in Volume~14. The race between the classes modulo $4$, drawn above, is a character sum; Chapter~\ref{ch:characters} defines the character, and Volume~6 proves Dirichlet's theorem.
\end{beauty}

\subsection*{Proof workshop: repair a false claim about Euclid's number}

Start with the true observation that $2\cdot3\cdot5+1=31$ is prime.
A few more computations might tempt you to claim that the product
of the first $k$ primes, plus $1$, is always prime. It is not:
\[
2\cdot3\cdot5\cdot7\cdot11\cdot13+1
=30031=59\cdot509.
\]
The calculation breaks the proposed formula for producing primes.
It does \emph{not} break Euclid's argument.

Take any finite collection $p_1,\ldots,p_k$ of primes and set
$N=p_1\cdots p_k+1$. The fundamental theorem of arithmetic gives
a prime $q$ dividing $N$. If $q$ were one of the $p_i$, it would
divide both $N$ and $N-1=p_1\cdots p_k$, hence would divide $1$.
No prime divides $1$. So $q$ is outside the list. The corrected
claim is both weaker and universal: every finite list misses a
prime. Notice exactly where factorisation, rather than primality
of $N$, enters the proof.


\section{Exercises}

\begin{exercise}[$\star$]\label{ex:mod4}
Show that there are infinitely many primes $p\equiv3\pmod4$. \emph{Hint: consider $4p_1\cdots p_k-1$.}
\end{exercise}

\begin{exercise}[$\star$]\label{ex:mersenne}
Show that if $2^n-1$ is prime then $n$ is prime, and that if $2^n+1$ is prime ($n\ge1$) then $n$ is a power of $2$.
\end{exercise}

\begin{exercise}[$\star$]\label{ex:germain}
Show that $n^4+4$ is composite for every integer $n>1$.
\end{exercise}

\begin{exercise}[$\star\star$]\label{ex:fermatnum}
Let $\mathcal F_n=2^{2^n}+1$ be the Fermat numbers. Show that $\mathcal F_0\mathcal F_1\cdots\mathcal F_{n-1}=\mathcal F_n-2$, deduce that the Fermat numbers are pairwise coprime, and give a new proof that there are infinitely many primes.
\end{exercise}

\begin{exercise}[$\star\star$]\label{ex:gaps}
Show that there are arbitrarily long runs of consecutive composite integers. Why does this not contradict Bertrand's postulate?
\end{exercise}

\begin{exercise}[$\star\star$]\label{ex:factorial}
Show that $n!$ is not a perfect power ($k$-th power with $k\ge2$) for $n\ge2$.
\end{exercise}

\begin{exercise}[$\star\star\star$]\label{ex:erdosdiv}
\emph{Erd\H{o}s' proof that $\sum\frac1p$ diverges.} Let $p_1<p_2<\cdots$ be the primes, and suppose $\sum_{j>k}\frac1{p_j}<\frac12$ for some $k$. Let $N\ge1$ be an integer.
\begin{enumerate}[label=(\alph*)]
\item Show that at most $\sum_{j>k}\frac N{p_j}<\frac N2$ integers $n\le N$ have a prime factor $p_j$ with $j>k$.
\item Writing $n=ab^2$ with $a$ squarefree, show that at most $2^k\sqrt N$ integers $n\le N$ have all their prime factors among $p_1,\dots,p_k$.
\item Conclude by taking $N=2^{2k+2}$.
\end{enumerate}
\end{exercise}

\begin{exercise}[$\star\star\star$]\label{ex:furstenberg}
\emph{Furstenberg's topological proof (1955).} Call $U\subset\Z$ open if it is a union of arithmetic progressions $a+b\Z$ with $b\ge1$.
\begin{enumerate}[label=(\alph*)]
\item Show that this defines a topology on $\Z$ in which every non-empty open set is infinite.
\item Show that each $a+b\Z$ is also closed.
\item If there were finitely many primes, show that $\{-1,1\}$ would be open, and conclude.
\end{enumerate}
\end{exercise}

\begin{exercise}[$\star$]\label{ex:t1-factorial-new-prime}
For $n\ge2$, show that every prime divisor of $n!+1$ exceeds $n$. Does it follow that $n!+1$ is prime? Settle the latter question using $n=5$.
\end{exercise}

\begin{exercise}[$\star\star$]\label{ex:t1-euclid-sequence}
Define $E_1=2$ and $E_{r+1}=E_1E_2\cdots E_r+1$. Compute $E_2,E_3,E_4,E_5$, and show that the $E_r$ are pairwise coprime. Must every $E_r$ be prime? Explain why this construction still proves the infinitude of primes.
\end{exercise}

\begin{exercise}[$\star$]\label{ex:t1-prod11}
Compute $\prod_{p\le 11}p$ and $\sum_{p\le 11}\ln p$, and compare them with $4^{11}$ and with $11\ln 4$.
\end{exercise}

\begin{beauty}
Euclid gives infinitely many primes, Euler shows they are plentiful (the reciprocal sum grows at least as fast as $\ln\ln x-1$), and Chebyshev pins their count to the order $\frac x{\ln x}$. The exact asymptotic, the prime number theorem, is the summit of Volume~6. Its classical proof uses complex analysis; the elementary proofs of Erd\H{o}s and Selberg show that complex analysis is not logically necessary. Primes in arithmetic progressions (Exercise~\ref{ex:mod4} is the easy case) need Dirichlet's characters, built in Volume~3 and put to work in Volume~6. And Exercise~\ref{ex:furstenberg} shows that even topology has something to say about primes.
\end{beauty}

\section{Hints}

\begin{hint}{ex:mod4}
Try to adapt Euclid's product-plus-one argument so that the new number is $3$ modulo $4$. If all its prime factors were $1$ modulo $4$, what would their product be modulo $4$? The construction $4p_1\cdots p_k-1$ gives the right contradiction.
\end{hint}
\begin{hint}{ex:mersenne}
Look at the exponent before looking at the huge number. If $n=ab$ with $a,b>1$, factor $(2^a)^b-1$. For the plus sign, an odd factor of the exponent yields the identity $X^{2r+1}+1=(X+1)(X^{2r}-X^{2r-1}+\cdots+1)$.
\end{hint}
\begin{hint}{ex:germain}
You do not need to guess two factors. Make the expression into a difference of squares: add and subtract $4n^2$ so that one piece becomes $(n^2+2)^2$. Then check that the two factors from $A^2-B^2$ are both greater than $1$.
\end{hint}
\begin{hint}{ex:fermatnum}
Multiplying expressions of the form $X-1$ and $X+1$ doubles the exponent. Use that observation inductively for consecutive Fermat numbers. The resulting product identity shows that any common divisor of two different Fermat numbers divides $2$. What does their parity tell you?
\end{hint}
\begin{hint}{ex:gaps}
Build the composite numbers instead of searching for an unusually large gap. The integers $(n+1)!+2,\ldots,(n+1)!+n+1$ each come with an obvious proper divisor. Bertrand's theorem concerns an interval whose length grows with its starting point, so compare those two scales.
\end{hint}
\begin{hint}{ex:factorial}
Bertrand's postulate gives a prime $p$ with $\frac n2<p\le n$. How many times does $p$ divide $n!$?
\end{hint}
\begin{hint}{ex:erdosdiv}
(a) The multiples of $p$ up to $N$ number $\lfloor N/p\rfloor$. (b) There are $2^k$ choices for $a$ and at most $\sqrt N$ for $b$. (c) The two counts must add up to at least $N$.
\end{hint}
\begin{hint}{ex:furstenberg}
(b) The complement of $a+b\Z$ is a union of $b-1$ progressions. (c) Zero is divisible by $2$, and every integer other than $0,\pm1$ has a prime factor, so $\Z\setminus\{-1,1\}=\bigcup_p p\Z$.
\end{hint}

\begin{hint}{ex:t1-factorial-new-prime}
If a prime $p\le n$ divides $n!+1$, then it divides both $n!$ and $n!+1$. For the counterexample, compute $5!+1$ rather than guessing from the first few values.
\end{hint}
\begin{hint}{ex:t1-euclid-sequence}
If $i<j$, reduce $E_j$ modulo $E_i$. The fifth term is $1807$; test divisibility by $13$. A new integer need not be prime to supply a new prime factor.
\end{hint}

\begin{hint}{ex:t1-prod11}
The primes are $2,3,5,7,11$. Theorem~\ref{thm:cheb}(i) is a strict inequality for the product.
\end{hint}

\section{Solutions}

\begin{solution}{ex:mod4}
\idea Euclid's construction needs a small adjustment when the new prime must lie in a particular residue class. Assume a finite list contains all primes congruent to $3$ modulo $4$. Form $4p_1\cdots p_k-1$. The new number is itself $3$ modulo $4$ and is coprime to every listed prime. Since a product of numbers congruent to $1$ modulo $4$ is still $1$ modulo $4$, at least one prime factor of the new number must lie in the missing class.
\writeup Suppose, toward a contradiction, that $p_1,\ldots,p_k$ list every prime congruent to $3$ modulo $4$. Construct
\[
N=4p_1p_2\cdots p_k-1.
\]
The number $N$ is odd and congruent to $3$ modulo $4$. Its prime factorisation must contain an odd number of factors congruent to $3$ modulo $4$, counted with multiplicity; otherwise the product would be $1$ modulo $4$. In particular, some prime $q\mid N$ satisfies $q\equiv3\pmod4$.

Could $q$ be one of the $p_i$? No: each $p_i$ divides $N+1=4p_1\cdots p_k$, so if it also divided $N$ it would divide $1$. Thus $q$ is a prime in the desired class that was missing from our supposedly complete list. This contradiction proves that there are infinitely many such primes.
\end{solution}

\begin{solution}{ex:mersenne}
\idea Do not begin by testing enormous numbers for primality. If an exponent factors, the corresponding expression often factors too. Put $n=ab$ into $2^n-1$ and use $X^b-1=(X-1)(X^{b-1}+\cdots+1)$. For a plus sign, check why an \emph{odd} quotient of exponents gives a similar factorisation. The arithmetic of the exponents decides which composite values cannot be prime.
\writeup First suppose that $2^n-1$ is prime. If $n=ab$ with $a,b>1$, then
\[
2^n-1=(2^a)^b-1
=(2^a-1)\bigl(1+2^a+\cdots+2^{a(b-1)}\bigr).
\]
Both factors exceed $1$, contradicting primality. Therefore $n$ must be prime.

Now suppose that $2^n+1$ is prime, with $n\ge1$. If $n$ had an odd factor $m>1$, write $n=km$. For odd $m$, the polynomial $X^m+1$ is divisible by $X+1$. Taking $X=2^k$ would give the proper factor $2^k+1$ of $2^n+1$. That is impossible. So $n$ has no odd factor greater than $1$, which means precisely that $n$ is a power of $2$.

These implications are necessary conditions, not converses: a prime exponent does not guarantee a Mersenne prime, and a power-of-two exponent does not guarantee a Fermat prime.
\end{solution}

\begin{solution}{ex:germain}
\idea Try to turn the expression into a difference of squares. Add and subtract $4n^2$: the expression $n^4+4n^2+4$ is a perfect square, and $4n^2$ is a perfect square too. Now use $A^2-B^2=(A-B)(A+B)$ and check that both factors are nontrivial for the allowed values of $n$. This is Sophie Germain's identity discovered by completing a square.
\writeup The expression does not look factorable until we complete a square:
\[
\begin{aligned}
n^4+4
&=(n^2+2)^2-(2n)^2\\
&=(n^2-2n+2)(n^2+2n+2).
\end{aligned}
\]
The first factor equals $(n-1)^2+1$, which is at least $2$ when $n>1$. The second is even larger. Both are proper positive factors of $n^4+4$, so the number is composite. The hypothesis $n>1$ is essential: for $n=1$ the number is $5$, a prime.
\end{solution}

\begin{solution}{ex:fermatnum}
\idea A telescoping product hides behind the numbers $2^{2^k}+1$. Multiply one such number by $2^{2^k}-1$ and use the difference of squares; the exponent doubles, so the result is the next expression minus $1$. This lets you express a Fermat number in terms of the earlier ones. If two Fermat numbers shared an odd prime factor, that prime would divide their difference, which is a power of $2$. Finish by checking parity.
\writeup Write $\mathcal F_n=2^{2^n}+1$. The product identity is proved by induction, using the difference of squares.

For $n=1$, $\mathcal F_0=3=\mathcal F_1-2$. If
$\prod_{j=0}^{n-1}\mathcal F_j=\mathcal F_n-2=2^{2^n}-1$, multiply by $\mathcal F_n=2^{2^n}+1$. The product becomes
\[
(2^{2^n}-1)(2^{2^n}+1)
=2^{2^{n+1}}-1=\mathcal F_{n+1}-2,
\]
which is the next case.

Now take $m<n$. A common divisor of $\mathcal F_m$ and $\mathcal F_n$ also divides $\prod_{j<n}\mathcal F_j=\mathcal F_n-2$, because that product includes $\mathcal F_m$. Hence it divides $2$. Every Fermat number is odd, so the common divisor is $1$. Thus the Fermat numbers are pairwise coprime.

Each $\mathcal F_n>1$ has a prime factor. Pairwise coprimality means that a prime factor chosen from one Fermat number cannot occur in any other. We have therefore found infinitely many distinct primes, without claiming that all the Fermat numbers themselves are prime.
\end{solution}

\begin{solution}{ex:gaps}
\idea To make a stretch of composite integers, give each position in the stretch its own guaranteed divisor. The factorial $(n+1)!$ is divisible by every integer $2,\ldots,n+1$. Therefore $(n+1)!+j$ is divisible by $j$ for each such $j$. Check that the resulting numbers are larger than their divisors, so they are genuinely composite. This construction proves arbitrarily long gaps without locating any particular next prime.
\writeup Fix a positive integer $n$ and set $M=(n+1)!$. Consider the $n$ consecutive integers
\[
M+2,\ M+3,\ \ldots,\ M+n+1.
\]
For each $j=2,\ldots,n+1$, the integer $j$ divides $M$ and hence divides $M+j$. Also $1<j<M+j$, so $j$ is a proper divisor. Every number in the displayed run is composite.

Why is Bertrand's theorem safe? The run has length $n$, but begins at a number of size roughly $(n+1)!$. Bertrand guarantees a prime between a number $m$ and $2m$, an interval of length $m$, not within every fixed number of steps. The composite run we constructed is tiny compared with its starting point. The two theorems measure gaps on different scales.
\end{solution}

\begin{solution}{ex:factorial}
\idea A perfect $k$-th power has every valuation divisible by $k$. So find one prime that divides $n!$ exactly once: a prime $p$ with $\frac n2<p\le n$ appears only in the factor $p$ itself. Bertrand guarantees such a prime.
\writeup Let $n\ge2$ and put $m=\lfloor n/2\rfloor$. Bertrand's theorem supplies a prime $p$ with $m<p\le2m\le n$. Since $p\ge m+1>n/2$, we have $p\le n<2p$, and therefore $p^2>n$.

Only one multiple of $p$ occurs among $1,2,\ldots,n$, namely $p$ itself, and no multiple of $p^2$ occurs. Legendre's formula gives
\[
v_p(n!)=\left\lfloor\frac np\right\rfloor=1.
\]
But if $n!=a^k$ with an integer $a$ and $k\ge2$, every prime valuation of $n!$ would be divisible by $k$. The valuation $1$ contradicts that. Thus $n!$ is not a perfect power.
\end{solution}

\begin{solution}{ex:erdosdiv}
\idea A counting argument by contradiction: if the tail of $\sum\frac1p$ were small, then few integers would have a ``large'' prime factor, so most integers up to $N$ would be built from $k$ small primes. But such integers are rare: writing $n=ab^2$ with $a$ squarefree, there are only $2^k$ choices for $a$ and $\sqrt N$ for $b$. Both counts cannot add up to $N$ when $N$ is large.
\writeup Assume, for a contradiction, that the series of prime reciprocals converges. Then for some $k$ its tail satisfies
\[
\sum_{j>k}\frac1{p_j}<\frac12.
\]

(a) At most $\lfloor N/p\rfloor\le N/p$ integers in $\{1,\ldots,N\}$ are divisible by a fixed prime $p$. By the union bound, fewer than $N/2$ of these integers have any prime factor among $p_{k+1},p_{k+2},\ldots$.

(b) Consider an integer whose prime factors all belong to the finite set $\{p_1,\ldots,p_k\}$. Write it uniquely as $n=ab^2$, with $a$ squarefree: put the even part of each prime exponent into $b^2$ and leave an exponent of $0$ or $1$ in $a$. There are at most $2^k$ choices for $a$, because it is a product of a subset of the first $k$ primes. Since $b^2\le n\le N$, there are at most $\lfloor\sqrt N\rfloor$ choices for $b$. Thus there are at most $2^k\sqrt N$ such integers.

(c) Every integer from $1$ to $N$ belongs to at least one of these two classes. We obtain
\[
N<\frac N2+2^k\sqrt N,
\qquad\text{hence}\qquad
\sqrt N<2^{k+1}.
\]
Choose $N=2^{2k+2}$. The final inequality would read $2^{k+1}<2^{k+1}$, a contradiction. The tail cannot be smaller than $1/2$, so the series $\sum_p1/p$ diverges.
\end{solution}

\begin{solution}{ex:furstenberg}
\idea (a), (b) are routine checks of the axioms. For (c), the key observation is that every integer except $\pm1$ is a multiple of some prime, so $\Z\setminus\{\pm1\}=\bigcup_pp\Z$. Each $p\Z$ is closed by (b); a \emph{finite} union of closed sets is closed, which would make the finite set $\{\pm1\}$ open, forbidden by (a).
\writeup (a) The empty set is an empty union of progressions, and $\Z=0+1\Z$ is a progression; both are open. Arbitrary unions are open by definition. For a finite intersection, suppose $x\in U\cap V$. Choose progressions $x+b\Z\subset U$ and $x+c\Z\subset V$. The smaller progression $x+bc\Z$ lies inside both, so every point of the intersection has an open progression neighbourhood. Thus the intersection is open. A nonempty open set contains a progression $a+b\Z$, which is infinite.

(b) The complement of a progression is a finite union of the other residue classes:
\[
\Z\setminus(a+b\Z)=\bigcup_{j=1}^{b-1}(a+j+b\Z).
\]
It is therefore open as well, so every progression is closed.

(c) Every integer other than $1$ and $-1$ has a prime divisor: this includes $0$, which every prime divides. Hence
\[
\Z\setminus\{-1,1\}=\bigcup_{\text{$p$ prime}}p\Z.
\]
If there were only finitely many primes, the right side would be a finite union of closed sets, hence closed. Its complement $\{-1,1\}$ would then be an open set. But this nonempty set is finite, contradicting part (a). There must be infinitely many primes.
\end{solution}

\begin{solution}{ex:t1-factorial-new-prime}
\idea The factorial contains every small prime, so adding $1$ forbids every one of them. A number can have only large prime factors and still have several of them.
\writeup Let $p$ be a prime factor of $n!+1$. If $p\le n$, then $p$ is one of the factors in $n!$, so $p\mid n!$. It would divide both $n!$ and $n!+1$, hence their difference $1$, which is impossible. Therefore every prime factor of $n!+1$ exceeds $n$.

This does not say the whole number is prime. For example,
\[
5!+1=121=11^2.
\]
The number is composite, but its only prime factor is $11>5$, exactly as the argument predicts.
\end{solution}

\begin{solution}{ex:t1-euclid-sequence}
\idea Each new term is one more than a multiple of every earlier term. A common divisor therefore has to divide $1$; the distinction between prime and pairwise coprime is the point of the experiment.
\writeup Compute successively from the definition:
\[
E_1=2,\quad E_2=3,\quad E_3=7,\quad E_4=43,\quad
E_5=2\cdot3\cdot7\cdot43+1=1807.
\]
The last term is already composite: $1807=13\cdot139$.

Now let $i<j$. By definition, $E_j-1$ is divisible by every earlier term, in particular $E_i$. So any common divisor of $E_i$ and $E_j$ also divides their difference $E_j-(E_j-1)=1$. We have proved that the terms are pairwise coprime.

Each $E_j>1$ has a prime factor. Since the terms are pairwise coprime, a prime chosen from one term cannot occur in another. There are therefore infinitely many distinct primes, even though the sequence itself contains composite numbers.
\end{solution}

\begin{solution}{ex:t1-prod11}
\idea Multiply the five primes, then take the natural logarithm.
\writeup The primes at most $11$ are $2,3,5,7,11$, and their product is
\[
2\cdot3\cdot5\cdot7\cdot11=2310.
\]
The bound from the chapter gives $2310<4^{11}=4194304$. Taking natural logarithms transforms the product into the sum
\[
\sum_{p\le11}\ln p=\ln2310\approx7.745
<11\ln4\approx15.249.
\]
Both estimates hold, as they must. This numerical check illustrates the inequalities but does not replace the proof for every $x$.
\end{solution}

\section*{Sources and further reading}

\begin{itemize}
\item G.~H.~Hardy and E.~M.~Wright, \emph{An Introduction to the Theory of Numbers}, 6th ed., Oxford University Press, 2008, Chapters~1--2 and~22. Euclid, Euler, Chebyshev and Bertrand, with historical notes.
\item M.~Aigner and G.~M.~Ziegler, \emph{Proofs from THE BOOK}, 6th ed., Springer, 2018, Chapters~1--2. Six proofs of the infinity of primes (including Furstenberg's and Erd\H{o}s' of Exercises~\ref{ex:furstenberg} and~\ref{ex:erdosdiv}) and Bertrand's postulate.
\item P.~Erd\H{o}s, Beweis eines Satzes von Tschebyschef, \emph{Acta Litt.\ Sci.\ Szeged} 5 (1932), 194--198. The original proof given in this chapter.
\item H.~Furstenberg, On the infinitude of primes, \emph{American Mathematical Monthly} 62 (1955), 353. The famous one-paragraph note.
\item R.~Crandall and C.~Pomerance, \emph{Prime Numbers: A Computational Perspective}, 2nd ed., Springer, 2005, Chapter~1. Sieves, $\pi(x)$ tables, and how primes are counted and found in practice.
\end{itemize}
